\chapter{AI Quiz Generation Engine}

\section{Pipeline Architecture and Operational Overview}
The module \texttt{recall/quiz\_engine.py} is the core of the system. It converts a request into a prompt, sends the prompt to an engine, and converts the reply into a validated quiz. The same pipeline serves document quizzes, topic quizzes and answer grading.

\section{Algorithmic Execution Flowchart}
\begin{figure}[H]
\centering
\begin{tikzpicture}[
    node distance=0.65cm,
    st/.style={draw, rounded corners, minimum width=5.4cm, minimum height=0.8cm, align=center, font=\small, fill=gray!8},
    dc/.style={draw, diamond, aspect=2.6, minimum width=3.8cm, align=center, font=\small, fill=yellow!10, inner sep=1pt},
    arr/.style={-{Stealth[length=2.2mm]}, thick}
]
\node[st] (p) {Build prompt (material, difficulty, counts, schema)};
\node[dc, below=of p] (g) {Use Groq?};
\node[st, below=of g] (call) {Call engine and obtain raw text};
\node[st, below=of call] (parse) {Strip code fences; extract outermost \{ \dots\ \}; \texttt{json.loads}};
\node[st, below=of parse] (val) {Validate: keep only well-formed questions};
\node[dc, below=of val] (any) {Any questions left?};
\node[st, below=of any] (ok) {Return quiz};
\node[st, right=2.3cm of call, fill=red!8] (fb) {Switch to Ollama\\ (after first failure)};
\node[st, right=2.3cm of any, fill=red!8] (err) {After 2 attempts:\\ raise error};

\draw[arr] (p) -- (g);
\draw[arr] (g) -- node[right, font=\scriptsize]{yes / auto} (call);
\draw[arr] (call) -- (parse);
\draw[arr] (parse) -- (val);
\draw[arr] (val) -- (any);
\draw[arr] (any) -- node[right, font=\scriptsize]{yes} (ok);
\draw[arr] (any.east) -- node[above, font=\scriptsize]{no} (err.west);
\draw[arr] (call.east) -- (fb.west);
\draw[arr] (fb.north) |- (p.east);
\end{tikzpicture}
\caption{Quiz generation and validation flow}
\label{fig:gen-flow}
\end{figure}

\section{Prompt Design}
Two system prompts are used. The document prompt instructs the model to act as a university examiner and to create questions \textit{only} from the provided material; the topic prompt allows well-established knowledge when no web material is available but forbids inventing facts. Both require a single JSON object with no markdown fences. The user prompt carries four ingredients:
\begin{enumerate}
\item the \textbf{difficulty guidance} string (Easy: recall and definitions; Medium: application with plausible distractors; Hard: multi-step analysis and edge cases);
\item the \textbf{question counts} for each type, with an instruction to return an empty list for any type whose count is zero;
\item the \textbf{JSON schema} (Listing~\ref{lst:quiz-json});
\item the \textbf{material}, truncated to 30{,}000 characters for the cloud engine or 10{,}000 for the local engine to respect its context window.
\end{enumerate}
The complete prompt text is reproduced in Appendix~\ref{app:prompts}.

\section{Parsing and Validation}
Models occasionally wrap JSON in prose or code fences. The parser removes fences, locates the first \texttt{\{} and the last \texttt{\}}, and decodes the slice. Validation then filters each category:
\begin{itemize}
\item an MCQ is kept only if it has exactly four options and an \texttt{answer\_index} from 0 to 3;
\item a written question needs both a question and a model answer;
\item a coding question needs both a question and a reference solution.
\end{itemize}
A reply that yields no valid questions is treated as a failure.

\section{Engines and Fallback Strategy}
\begin{table}[H]
\centering
\begin{tabularx}{\textwidth}{@{} p{2.4cm} p{4.4cm} X @{}}
\toprule
\textbf{Engine} & \textbf{Selection} & \textbf{Details} \\
\midrule
Groq Cloud & \texttt{engine=groq}, or \texttt{auto} when \texttt{GROQ\_API\_KEY} starts with \texttt{gsk\_} & HTTPS call to the chat-completions endpoint with \texttt{response\_format=json\_object}, temperature 0.3, 45~s time-out. \\
Ollama & \texttt{engine=ollama}, or \texttt{auto} without a key & Local model (default \texttt{llama3.2:3b}) called through the \texttt{ollama} package in JSON mode. \\
\bottomrule
\end{tabularx}
\caption{Supported engines}
\label{tab:engines}
\end{table}
The function \texttt{\_run\_prompt} makes up to two attempts. If the first attempt with the cloud engine fails for any reason, the second attempt uses the local engine. After two failures a descriptive error is raised, which the server turns into HTTP~502.

\section{AI Grading of Written Answers}
\texttt{grade\_short\_answer} sends the model answer, the key points and the student's answer to the engine and asks for a JSON object with \texttt{score} (0--100), \texttt{hits}, \texttt{misses} and \texttt{feedback}, using temperature 0.2 for consistency. If grading fails, a neutral result (score~0, all key points listed as missed, and a message asking the learner to compare with the model answer) is returned so that the quiz can continue.

\section{Markdown Export}
\texttt{to\_markdown} renders a quiz as Markdown: numbered MCQs with lettered options, the answer line with explanation and optional source quote, model answers for written questions and reference solutions in fenced code blocks.

\screenshot{loading_generation}{Generation in progress: the loading view shows the stages ``Searching the web'', ``Reading the best sources'', ``Writing your questions'' and ``Checking answers''}{fig:loading}{Click Generate on the home page and capture the animated loading screen with its step list.}

\screenshot{ai_grading}{AI grading of a written answer: score, overall feedback, key points covered (green) and missed (red)}{fig:ai-grading}{Answer a written question and press ``Check with AI''.}
